\section*{Chapitre \ref{ch:b2:potential-wells-tunneling} --- Puits de potentiel, barrières et effet tunnel}

\begin{solution}{exo:b2:potential-wells-tunneling:1}
\qty{1}{nm}~: \qty{0.38}{}, \qty{1.50}{}, \qty{3.38}{eV}~; $2 \to 1$~: \qty{1.13}{eV}, \qty{1.1}{\micro m}.
\qty{0.1}{nm}~: \qty{38}{eV}. Nucléon~: \qty{8}{MeV}. Bille~: $E_1 = \qty{5.5e-63}{J}$~;
$n = 2mvL/h = 3 \times 10^{27}$.
\end{solution}

\begin{solution}{exo:b2:potential-wells-tunneling:2}
Confinement $\qty{3.76}{eV.nm^2}/L^2$~: $0.10$, $0.24$, $0.42$, \qty{0.60}{eV}~; total
$1.84$, $1.98$, $2.16$, \qty{2.34}{eV}~: \qty{673}{}, \qty{628}{}, \qty{574}{}, \qty{530}{nm} ---
rouge, orange, jaune-vert, vert.
\end{solution}

\begin{solution}{exo:b2:potential-wells-tunneling:3}
$\kappa = \qty{1.09e10}{m^{-1}}$~; $\eu^{-2\kappa a}$~: $1.5 \times 10^{-3}$, $1.8 \times 10^{-5}$,
$3 \times 10^{-10}$, $10^{-19}$~; $\times 8.8$ par \qty{0.1}{nm}~; proton~: $\kappa$ est
$43$ fois plus grand, $\eu^{-280}$ à \qty{0.3}{nm}.
\end{solution}

\begin{solution}{exo:b2:potential-wells-tunneling:4}
$k_1 = \qty{7.3e9}{m^{-1}}$, $k_2 = \qty{5.1e9}{m^{-1}}$~; $R = 0.03$, $T = 0.97$. Marche
descendante~: $k_2 = \qty{8.9e9}{m^{-1}}$, $R = 0.01$. Une particule classique ne
rebrousse jamais chemin quand elle en a l'énergie~; une onde lumineuse est partiellement réfléchie
à tout changement d'indice.
\end{solution}

\begin{solution}{exo:b2:potential-wells-tunneling:5}
(a), (b) \cref{prop:b2:potential-wells-tunneling:infinite}. (c) $L/2$~;
$0.18L$~; $0$~; $\pi\hbar/L$~; produit $0.57\hbar > \hbar/2$. (d) Des oscillations
rapides autour de la moyenne $1/L$ --- la densité uniforme classique.
\end{solution}

\begin{solution}{exo:b2:potential-wells-tunneling:6}
(a) \cref{prop:b2:potential-wells-tunneling:finite}. (b) $1$, $3$, $7$.
(c) $R = 2.6$, deux états~; $\kappa = \qty{4.4e9}{m^{-1}}$, profondeur \qty{0.23}{nm}. (d)
Une infinité, tendant par en dessous vers celles du puits infini.
\end{solution}

\begin{solution}{exo:b2:potential-wells-tunneling:7}
(a) $1 + r = t$, $\iu k(1 - r) = -\kappa t$. (b) Numérateur et dénominateur
sont conjugués~; phase $-2\arctan(\kappa/k)$. (c) $\psi^*\psi'$ est réel pour une
exponentielle réelle. (d) \qty{0.2}{nm}~; comme l'\omterm{def:b2:dispersion-wave-packets:complex}{onde évanescente} de la \omterm{thm:b2:wave-interfaces:snell}{réflexion
totale}~: présente, ne portant aucun flux, et capable d'alimenter un second milieu
--- c'est cela, l'\omterm{thm:b2:potential-wells-tunneling:tunnel}{effet tunnel}.
\end{solution}

\begin{solution}{exo:b2:potential-wells-tunneling:8}
(a) \cref{thm:b2:potential-wells-tunneling:tunnel}. (b) $\kappa a = 3$~:
exact $1/(1 + \sinh^23) = 9.9 \times 10^{-3}$, approximation $4\eu^{-6} = 9.9 \times
10^{-3}$~; $\kappa a = 1$~: $0.42$ contre $0.54$. (c) $\kappa \to \iu k_2$~; $T = 1$
pour $k_2a = n\pi$. (d) Les réflexions des deux faces s'annulent quand la
barrière fait un nombre entier de demi-longueurs d'onde --- transmission
résonnante.
\end{solution}

\begin{solution}{exo:b2:potential-wells-tunneling:9}
(a) $b = \qty{52}{fm}$~; $G = 82$~; $T = \eu^{-82} \approx 2 \times 10^{-36}$. (b) $v =
\qty{1.55e7}{m/s}$, $10^{21}$ chocs par seconde~: taux $2 \times 10^{-15}\,\unit{s^{-1}}$,
durée de vie $\sim 10^7$ ans. (c) $G = 100$ et $69$~: $\eu^{18} \approx 10^8$ fois plus longue,
$\eu^{-13} \approx 10^{-6}$ fois plus courte. (d) L'$\alpha$ est fortement liée et part
avec une énergie positive~; un proton ne le ferait pas~; un $^{12}$C a trois fois la
charge et un exposant bien plus grand (la radioactivité de grappe existe, à $10^{-10}$).
\end{solution}

\begin{solution}{exo:b2:potential-wells-tunneling:10}
(a) \qty{24}{GHz}~; \qty{41}{ps}. (b) $\psi = (\varphi_+\eu^{-\iu E_+t/\hbar} +
\varphi_-\eu^{-\iu E_-t/\hbar})/\sqrt2$~; $P = \sin^2(\pi\nu t)$. (c) $\eu^{(\sqrt2 - 1)
\kappa a} = 15$~: $\kappa a = 6.5$. (d) Des groupes plus lourds et des barrières plus hautes rendent
la période tunnel plus longue que des années.
\end{solution}

\begin{solution}{exo:b2:potential-wells-tunneling:11}
(a) \qty{1.09e10}{m^{-1}}~; $\times 8.8$~; $1\%$ de $I$ font \qty{0.5}{pm}. (b) $6 \times 10^9$.
(c) Chaque voisin à $d' = \qty{0.56}{nm}$ porte $\eu^{-2\kappa \times 0.06\,\text{nm}}
= 0.28$ du courant de l'apex~; l'apex en porte environ la moitié. (d)
$10^{-5}\,\unit{K}$ --- impossible~; des architectures symétriques où pointe et échantillon
se dilatent ensemble, des balayages rapides, des matériaux à faible dilatation, et l'asservissement.
\end{solution}

\begin{solution}{exo:b2:potential-wells-tunneling:12}
(a) \qty{0.056}{} et \qty{0.008}{eV}. (b) \qty{1.48}{eV}, \qty{835}{nm}. (c) Confinement
\qty{0.17}{eV}~: $L = \qty{6.1}{nm}$. (d) $R = 3.6$~: trois.
\end{solution}

\begin{solution}{pb:b2:potential-wells-tunneling:1}
\textbf{1.} $\kappa = \sqrt{2m\Phi}/\hbar = \qty{1.09e10}{m^{-1}}$.

\textbf{2.} $\eu^{2.2} = 8.8$~; $\eu^{0.022}$~: $2.2\%$ par picomètre.

\textbf{3.} \qty{0.5}{pm}.

\textbf{4.} $6 \times 10^9$ par seconde~; \qty{0.16}{ns} d'écart, contre $\hbar/\Phi
\sim 10^{-16}\,\unit{s}$~: un à la fois.

\textbf{5.} Environ la moitié~; le reste décroît si vite avec la distance que
l'image est celle d'un seul atome.

\textbf{6.} $\times 8.8$~; le courant constant garde la pointe hors de danger et change
l'exponentielle en une hauteur linéaire.

\textbf{7.} $\Delta T = 10^{-12}/(10^{-5} \times 0.02) = \qty{5e-6}{K}$~; des montages
symétriques, la rapidité, des matériaux à faible dilatation, et un asservissement qui suit la
dérive.

\textbf{8.} Un trapèze, abaissé du côté collecteur~: la barrière
effective s'amincit et $I$ croît plus vite que linéairement.

\textbf{9.} $\kappa' = \qty{0.89e10}{m^{-1}}$~; à \qty{0.5}{nm} l'exposant tombe de
$10.9$ à $8.9$, $I \times 7$~: la pointe recule de $\ln 7/2\kappa' \approx \qty{0.1}{nm}$
--- une bosse électronique, non géométrique.

\textbf{10.} $V(R_{\text{n}}) = 2 \times 90 \times 1.44/8 = \qty{32}{MeV}$~; $b = \qty{62}{fm}$.

\textbf{11.} $v = \qty{1.4e7}{m/s}$~; $9 \times 10^{20}$ chocs par seconde.

\textbf{12.} $G = 96$~; $T = \eu^{-96} \approx 10^{-42}$.

\textbf{13.} Taux $\sim 10^{-21}\,\unit{s^{-1}}$, demi-vie $\sim 10^{13}$ ans --- mille
fois la valeur mesurée~: pour un exposant de cent, un modèle
grossier s'en tire bien en tombant à quelques ordres de grandeur près.

\textbf{14.} $G = 79$~: $\eu^{17} \approx 10^7$ fois plus courte, environ $10^6$ ans ---
un MeV de plus, sept ordres de grandeur.

\textbf{15.} $G = 92$~: $\eu^{4.5} \approx 90$ fois plus courte --- un fermi de rayon
vaut deux ordres de grandeur.

\textbf{16.} Un proton n'est pas préformé avec une énergie positive (il est
lié par quelque \qty{7}{MeV})~; un $^{12}$C a trois fois la charge et un
exposant trois fois plus grand.

\textbf{17.} Oui~: l'$\alpha$ a toujours \qty{4.2}{MeV}~; sa \omterm{def:b2:schrodinger-wave-functions:psi}{fonction d'onde}
a seulement une partie évanescente sous la barrière, où aucune mesure
ne la trouve avec une énergie cinétique négative.

\textbf{18.} \qty{23.7}{GHz}~; \qty{1.27}{cm}.

\textbf{19.} $(|{\uparrow}\rangle \pm |{\downarrow}\rangle)/\sqrt2$~; la
symétrique est la plus basse --- pas de nœud, moins de courbure.

\textbf{20.} $\psi = (\varphi_+\eu^{-\iu E_+t/\hbar} + \varphi_-\eu^{-\iu E_-t/\hbar})/\sqrt2$~;
$P_{\downarrow} = \sin^2(\pi\nu t)$~; période \qty{42}{ps}.

\textbf{21.} $\eu^{-\Delta/k_BT} = 0.996$~: presque égales, pas d'inversion~; un
champ électrique inhomogène trie le jet et envoie les molécules de l'état
supérieur dans la cavité.

\textbf{22.} \qty{1.6e-23}{J}~; $6 \times 10^{12}$ molécules par seconde.

\textbf{23.} $\kappa a = 6.5$.

\textbf{24.} Elle est fixée par la molécule seule, identique pour chaque
molécule, et insensible à l'extérieur~; les décalages viennent de
l'\omterm{thm:b2:sound-waves:doppler}{effet Doppler} dans le jet, des champs électriques parasites, des collisions et
de l'entraînement par la cavité.

\textbf{25.} Le confinement quantifie ($E_n \propto n^2/L^2$)~; l'évanescence laisse
passer l'onde ($T \sim \eu^{-2\kappa a}$)~; les exposants sont fixés par
$\sqrt{2m(V_0 - E)}/\hbar$ et la largeur --- masse, hauteur, distance.
\end{solution}
